%============================================================
% CHƯƠNG 6: KẾT LUẬN VÀ HƯỚNG PHÁT TRIỂN
%============================================================
\chapter{KẾT LUẬN VÀ HƯỚNG PHÁT TRIỂN}
\label{chap:conclusion}

Chương này tổng kết quá trình thực hiện đồ án, đánh giá những thành
tựu đạt được so với mục tiêu ban đầu và nhận diện các hạn chế còn
tồn tại. Cuối cùng, chương đề xuất các hướng phát triển tiềm năng
để nâng cao chất lượng và mở rộng phạm vi ứng dụng của hệ thống
trong tương lai.

%------------------------------------------------------------
\section{Kết luận}
\label{sec:conclusion:summary}
%------------------------------------------------------------

Trước thực trạng các hệ thống quan trắc môi trường hiện tại hoặc có
chi phí quá cao, hoặc thiếu khả năng dự báo và trực quan hóa không
gian, đồ án đã đề xuất và triển khai thành công một \textbf{hệ thống
AIoT giám sát và dự báo chất lượng không khí} hoàn chỉnh với chi phí
thấp và mã nguồn mở.

\subsection{Kết quả đạt được}
\label{subsec:conclusion:result}

Về mặt kỹ thuật, hệ thống đã đạt được các thành tựu cụ thể sau:

\begin{itemize}[leftmargin=2cm]
  \item \textbf{Node cảm biến IoT:} Các node ESP32/DHT11 hoạt động
    ổn định, thu thập và gửi dữ liệu định kỳ 5 giây qua MQTT với
    cơ chế tự động đăng ký (provisioning) và khôi phục kết nối
    (reconnect) không cần can thiệp thủ công.

  \item \textbf{Backend xử lý thời gian thực:} Kiến trúc Node.js
    tích hợp Aedes MQTT broker và Socket.io hoạt động ổn định,
    phân phối dữ liệu đến dashboard với độ trễ end-to-end trung
    bình dưới 500ms trong điều kiện mạng LAN.

  \item \textbf{Pipeline dự báo LSTM:} Mô hình LSTM (PyTorch) với
    window size 30 và forecast horizon 15 phút được huấn luyện hội
    tụ tốt, phục vụ inference qua FastAPI microservice trong dưới
    200ms mỗi lần dự báo.

  \item \textbf{Nội suy không gian Kriging/IDW:} Từ $n$ điểm đo
    rời rạc, hệ thống tái tạo bản đồ phân bố nhiệt độ, độ ẩm và AQI
    liên tục trên lưới 40$\times$50 điểm, render heatmap trực quan
    trên Leaflet trong dưới 2 giây.

  \item \textbf{Dashboard React:} Giao diện người dùng hiển thị dữ
    liệu realtime, biểu đồ lịch sử, kết quả dự báo và heatmap
    không gian trên cùng một nền tảng web, không cần tải lại trang.

  \item \textbf{Chi phí thấp:} Tổng chi phí phần cứng một node
    cảm biến dưới 150.000 VND, thấp hơn nhiều lần so với các
    thiết bị thương mại tương đương, phù hợp triển khai mật độ cao.
\end{itemize}

\subsection{Đánh giá so với mục tiêu}
\label{subsec:conclusion:eval}

So sánh với 6 mục tiêu cụ thể đề ra ở Chương~\ref{chap:intro},
hệ thống đã đáp ứng đầy đủ tất cả các mục tiêu trong phạm vi
prototype. Đặc biệt, sự kết hợp giữa \textbf{dự báo chuỗi thời
gian} (LSTM) và \textbf{nội suy không gian} (Kriging/IDW) trong
cùng một hệ thống là điểm kỹ thuật nổi bật, chưa được tích hợp
trong các giải pháp DIY tương tự hiện có.

\subsection{Hạn chế}
\label{subsec:conclusion:limit}

Tuy nhiên, đồ án vẫn còn một số hạn chế cần nhìn nhận thẳng thắn:

\begin{itemize}[leftmargin=2cm]
  \item \textbf{Cảm biến hạn chế:} Hệ thống hiện chỉ đo nhiệt độ
    và độ ẩm qua DHT11, chưa tích hợp các cảm biến quan trọng hơn
    như PM2.5, CO$_2$ hay VOC vốn liên quan trực tiếp đến chỉ số
    AQI.

  \item \textbf{Độ chính xác DHT11:} Cảm biến DHT11 có sai số
    $\pm 2^\circ$C và $\pm 5\%$RH, chưa đạt chuẩn đo lường y tế
    hay môi trường chính thức.

  \item \textbf{Dữ liệu huấn luyện hạn chế:} Mô hình LSTM được
    huấn luyện trên tập dữ liệu thu thập trong môi trường thực
    nghiệm quy mô nhỏ, chưa đủ đa dạng để đảm bảo tổng quát hóa
    tốt trong mọi điều kiện thực tế.

  \item \textbf{Nội suy không gian với ít node:} Kriging cho kết
    quả tốt nhất khi có nhiều điểm đo. Với chỉ 2--3 node, kết quả
    nội suy còn phụ thuộc nhiều vào vị trí đặt node và có thể chưa
    phản ánh chính xác thực tế.

  \item \textbf{Chưa có đánh giá lâm sàng:} Các kết quả thực
    nghiệm chỉ mang tính tham khảo trong điều kiện phòng lab,
    chưa được kiểm chứng theo tiêu chuẩn đo lường quốc gia hay
    so sánh với thiết bị tham chiếu chuẩn.

  \item \textbf{Bảo mật chưa hoàn thiện:} Hệ thống chưa mã hóa
    kênh truyền MQTT (chưa dùng MQTT over TLS), tiềm ẩn rủi ro
    bảo mật khi triển khai trên mạng công cộng.

  \item \textbf{Thông tin kết nối cơ sở dữ liệu:} Trong giai đoạn
    đầu phát triển, thông tin xác thực PostgreSQL
    (\texttt{DB\_PASSWORD}) được hardcode trực tiếp trong các file
    Python ML (\texttt{forecast/*.py}). Hệ thống đã được khắc phục
    bằng cách chuyển toàn bộ sang đọc từ biến môi trường thông qua
    file \texttt{.env} (sử dụng thư viện \texttt{python-dotenv}),
    phù hợp với thực hành bảo mật tốt nhất (security best
    practices) trong phát triển phần mềm.
\end{itemize}

%------------------------------------------------------------
\section{Hướng phát triển}
\label{sec:conclusion:future}
%------------------------------------------------------------

Dựa trên những kết quả đạt được và hạn chế đã nhận diện, đồ án
đề xuất các hướng phát triển cụ thể như sau:

\subsection{Mở rộng khả năng cảm biến}
\label{subsec:future:sensor}

Ưu tiên hàng đầu là tích hợp thêm các cảm biến quan trọng để đo
được chỉ số AQI thực sự: cảm biến PMS5003 hoặc SDS011 đo PM2.5
và PM10, cảm biến MH-Z19 đo CO$_2$, và cảm biến SGP30 đo VOC.
Thay thế DHT11 bằng DHT22 hoặc SHT31 để cải thiện độ chính xác.
Việc bổ sung cảm biến GPS cho phép triển khai ngoài trời với
định vị tự động thay vì nhập tọa độ thủ công.

\subsection{Nâng cao mô hình Machine Learning}
\label{subsec:future:ml}

Nghiên cứu và so sánh các kiến trúc mô hình tiên tiến hơn như
Temporal Fusion Transformer (TFT) hoặc N-BEATS cho bài toán dự báo
đa biến, đặc biệt khi số lượng features tăng lên. Áp dụng
\textbf{federated learning} để huấn luyện mô hình trên dữ liệu từ
nhiều node phân tán mà không cần tập trung dữ liệu về máy chủ,
đảm bảo quyền riêng tư. Xây dựng pipeline tự động tái huấn luyện
(auto-retraining) khi phát hiện model drift.

\subsection{Cải thiện nội suy không gian}
\label{subsec:future:spatial}

Nghiên cứu các phương pháp nội suy không gian nâng cao như
Co-Kriging (kết hợp nhiều biến tương quan) hoặc tích hợp mô hình
CFD (Computational Fluid Dynamics) cho các không gian phức tạp.
Xây dựng cơ chế tự động xác định vị trí tối ưu để đặt thêm node
cảm biến dựa trên phân tích phương sai Kriging.

\subsection{Tối ưu bảo mật và hạ tầng}
\label{subsec:future:infra}

Triển khai \textbf{MQTT over TLS} để mã hóa toàn bộ kênh truyền
từ node IoT đến broker. Áp dụng mô hình \textbf{zero-trust} với
xác thực thiết bị bằng certificate thay vì chỉ dùng api\_key.
Containerize toàn bộ hệ thống bằng Docker Compose để dễ dàng
triển khai trên nhiều môi trường. Đánh giá khả năng mở rộng
(scalability) với hàng trăm node đồng thời thông qua load testing.

\subsection{Tích hợp hệ sinh thái và mở rộng ứng dụng}
\label{subsec:future:integration}

Phát triển ứng dụng di động (React Native) để người dùng có thể
theo dõi và nhận cảnh báo mọi lúc mọi nơi. Tích hợp với các nền
tảng smarthome phổ biến như Home Assistant hoặc Google Home để
kết hợp điều khiển thiết bị dựa trên dữ liệu môi trường (ví dụ:
tự động bật điều hòa khi nhiệt độ vượt ngưỡng). Xây dựng
\textbf{public API} và tài liệu hướng dẫn để cộng đồng nghiên
cứu có thể tích hợp dữ liệu vào các ứng dụng khác.

\subsection{Thử nghiệm và đánh giá quy mô lớn}
\label{subsec:future:evaluation}

Tiến hành triển khai thử nghiệm thực tế với 10--20 node trong môi
trường thực (văn phòng, trường học, khu dân cư) và so sánh kết
quả với thiết bị đo chuẩn để đánh giá độ chính xác khách quan.
Thực hiện nghiên cứu người dùng (user study) để đánh giá tính
khả dụng của dashboard và cải thiện UX dựa trên phản hồi thực tế.
Công bố tập dữ liệu thu thập được và mã nguồn hệ thống như một
đóng góp mở cho cộng đồng nghiên cứu AIoT và quan trắc môi trường.
